\documentclass[12pt, a4paper]{article}

\usepackage[utf8]{inputenc}
\usepackage[T1]{fontenc}
\usepackage{lmodern}
\usepackage[english]{babel}
\usepackage{geometry}
\usepackage{titlesec}
\usepackage{enumitem}
\usepackage{hyperref}
\usepackage{parskip}
\usepackage{booktabs}
\usepackage{array}
\usepackage{xcolor}
\usepackage{fancyhdr}
\usepackage{graphicx}
\usepackage{microtype}

\geometry{
  top=2.5cm, bottom=2.5cm,
  left=3cm,  right=2.5cm
}

\hypersetup{
  colorlinks=true,
  linkcolor=black,
  urlcolor=blue,
  pdftitle={Activity Report --- Vezeteu Andrei},
}

\pagestyle{fancy}
\fancyhf{}
\rhead{\small Activity Report}
\lhead{\small Vezeteu Andrei}
\cfoot{\thepage}
\renewcommand{\headrulewidth}{0.4pt}

\titleformat{\section}{\large\bfseries}{}{0em}{}[\titlerule]
\titleformat{\subsection}{\normalsize\bfseries}{}{0em}{}

\setlist[itemize]{noitemsep, topsep=4pt}

% ---------------------------------------------------------------
\begin{document}

% ===== HEADER ===================================================
\begin{center}
  {\large \textbf{``Alexandru Ioan Cuza'' University of Iasi}} \\[4pt]
  {\normalsize Faculty of Computer Science} \\[16pt]

  \rule{\linewidth}{0.5pt} \\[10pt]

  {\LARGE \textbf{Activity Report}} \\[6pt]
  {\large Bachelor's Thesis} \\[16pt]

  \rule{\linewidth}{0.5pt}
\end{center}

\vspace{10pt}

\begin{tabular}{@{}ll}
  \textbf{Title:}            & \textit{Sequre: A Real-Time Simulation and Visualization} \\
                             & \textit{Platform for the BB84 Quantum Key Distribution Protocol} \\[4pt]
  \textbf{Student:}          & Vezeteu Andrei \\[4pt]
  \textbf{Supervisor:}       & Conf.\ Dr.\ Andreea Arusoaie \\[4pt]
  \textbf{Academic year:}    & 2025 -- 2026 \\
\end{tabular}

\vspace{16pt}

% ===== 1. INTRODUCTION ==========================================
\section{Introduction}

This thesis aims to design and implement an interactive web platform for the
simulation and visualization of the \textbf{BB84} protocol --- the first
\textit{Quantum Key Distribution} (QKD) scheme, proposed by Charles Bennett
and Gilles Brassard in 1984. The protocol enables two parties, Alice and Bob,
to establish a shared secret key over an insecure channel, with security
guaranteed by the laws of quantum mechanics rather than computational hardness.

The platform, named \textbf{Sequre}, provides a complete, end-to-end
demonstration of the quantum cryptographic pipeline: qubit generation and
transmission, basis sifting, Quantum Bit Error Rate (QBER) estimation, error
correction, and privacy amplification. A distinguishing feature of the project
is its real integration with \textbf{IBM Quantum} hardware, allowing simulations
to run on both classical noise simulators and physical quantum processors.
      
Beyond its pedagogical role, Sequre is conceived as a research-oriented
testbed: its broader goal is to \textbf{probe the practical limits of
quantum key distribution and quantum cryptography} under realistic noise,
hardware imperfections, and adversarial conditions. The long-term ambition
is for the platform to evolve from an academic prototype into a tool with
\textbf{real-world applicability} --- usable by researchers, educators, and,
eventually, as a reference implementation for the integration of QKD
primitives into production cryptographic systems.

% ===== 2. WORK CARRIED OUT =====================================
\section{Work Carried Out}

\subsection{2.1 System Architecture}

The application follows a client-server architecture:

\begin{itemize}
  \item \textbf{Backend:} FastAPI (Python), providing real-time communication
        via WebSockets and persistent storage using SQLite.
  \item \textbf{Frontend:} React (Vite) with Tailwind CSS, consuming a
        WebSocket stream to animate simulation events live in the browser.
  \item \textbf{Runtime environment:} Docker Compose for service isolation
        and reproducibility across development environments.
\end{itemize}

\subsection{2.2 BB84 Protocol Implementation}

The full BB84 protocol was implemented using \textbf{Qiskit} for quantum
circuit construction and \textbf{Qiskit-Aer} for noisy simulation:

\begin{itemize}
  \item Alice--Bob quantum circuits with configurable basis choices (Z and X).
  \item Composite noise model: depolarizing errors and measurement errors,
        with user-configurable probabilities.
  \item Basis sifting: retention of bits where Alice and Bob chose the same
        measurement basis (~50\% of raw qubits).
  \item QBER calculation and security decision using an 11\% abort threshold.
  \item Error correction via 4-bit parity blocks (simplified CASCADE-style).
  \item Privacy amplification via SHA-256 truncated to a 128-bit final key.
\end{itemize}

\subsection{2.3 Eve Eavesdropping Simulation}

An \textit{intercept-resend} attack model was implemented with two modes:

\begin{itemize}
  \item \textbf{Weak Eve:} intercepts 30\% of qubits, producing a moderate
        QBER increase that may remain below the security threshold.
  \item \textbf{Strong Eve:} intercepts 100\% of qubits, generating a QBER
        of approximately 25\%, well above the 11\% abort threshold.
\end{itemize}

Both modes demonstrate the no-cloning theorem and Heisenberg uncertainty
principle in practice: Eve's forced measurement inevitably disturbs the
quantum states, leaving a detectable trace in the error rate.

\subsection{2.4 IBM Quantum Hardware Integration}

The platform supports executing simulations on real quantum hardware via
\textbf{Qiskit IBM Runtime} (\textit{SamplerV2}), with support for the
\texttt{ibm\_fez}, \texttt{ibm\_marrakesh}, and \texttt{ibm\_kingston}
backends. Circuit transpilation and job polling are handled asynchronously,
with live status messages streamed to the frontend via WebSocket.

\subsection{2.5 User Interface}

The following visualization components were implemented:

\begin{itemize}
  \item \textbf{PhotonGrid:} cell-by-cell animated visualization of the
        first 80 qubits exchanged. Eve-intercepted qubits are highlighted
        with an orange border and a badge indicator. Each cell encodes
        Alice's basis symbol, bit value, and Bob's basis.
  \item \textbf{SimulationSummary:} a full key distillation pipeline
        (Raw $\to$ Sifted $\to$ After EC $\to$ Final Key, 128-bit),
        a contextual security status banner, per-run statistics, and
        a final secret key display with a one-click copy button.
  \item \textbf{QBERChart:} a line chart of QBER values across all
        previous simulation runs, with the 11\% security threshold
        marked as a reference line.
  \item \textbf{EducationalPanel:} five accordion cards covering the
        core concepts of quantum cryptography: the BB84 protocol, basis
        sifting, QBER interpretation, the no-cloning theorem and
        Heisenberg uncertainty principle, and error correction with
        privacy amplification.
\end{itemize}

\subsection{2.6 Data Persistence}

Each completed simulation is automatically persisted to a SQLite database,
storing configuration parameters, QBER, sifted key length, and security
status. REST endpoints are exposed for history retrieval and deletion
(\texttt{GET /api/v1/history}, \texttt{DELETE /api/v1/history}).

% ===== 3. CURRENT STATUS ========================================
\section{Current Status of the Project}

At the time of writing this report, all major components of the platform
are \textbf{fully functional and integrated}. The simulator produces results
consistent with the theoretical predictions of the BB84 protocol: the sifting
rate is approximately 50\%, Strong Eve reliably generates QBER $\approx$ 25\%,
and Weak Eve produces values below the security threshold in roughly 70\% of
runs, consistent with its 30\% interception rate.

The source code is version-controlled on the \texttt{main} branch, with a
documented commit history on GitHub.

% ===== 4. FUTURE WORK ===========================================
\section{Future Work}

The following extensions are identified as meaningful directions for future
development, grounded in current research in quantum cryptography:

\begin{itemize}

  \item \textbf{Decoy-state BB84 with weak coherent pulse modelling.}
        Practical QKD systems do not use ideal single-photon sources; instead,
        they rely on attenuated laser pulses (weak coherent pulses, WCP) that
        occasionally emit multi-photon states. This makes them vulnerable to
        the \textit{Photon Number Splitting} (PNS) attack, in which an
        eavesdropper retains extra photons without disturbing the transmission.
        The decoy-state method~\cite{lo2005decoy} mitigates this by randomly
        injecting pulses of varying intensities, allowing Alice and Bob to
        estimate single-photon channel statistics and detect PNS attempts.
        Incorporating a WCP source model and decoy-state parameter
        optimisation into Sequre would substantially improve the realism of
        the simulation~\cite{arxiv2405}.

  \item \textbf{Channel loss and key rate vs.\ distance modelling.}
        The current platform models gate-level noise (depolarising, measurement
        error), but does not account for \emph{photon loss} due to transmission
        distance over optical fibre or free-space links. Adding an
        attenuation model (loss coefficient $\alpha$ in dB/km for fibre, or
        beam-divergence loss for free-space) and plotting the resulting
        \emph{secure key rate as a function of channel length} would bring
        the simulator in line with standard QKD benchmarking
        methodology~\cite{arxiv2505, arxiv2503}.

  \item \textbf{CASCADE error correction algorithm.}
        The current scheme discards entire 4-bit parity blocks on mismatch,
        leading to significant key loss at higher QBER values. The CASCADE
        algorithm~\cite{brassard1994} performs iterative multi-pass
        reconciliation, identifying and correcting individual erroneous bits
        rather than whole blocks, and achieves near-optimal information
        reconciliation efficiency. Replacing the existing scheme with CASCADE
        would make post-processing suitable for realistic noisy channels.

  \item \textbf{Support for additional QKD protocols.}
        Implementing \textbf{E91}~\cite{ekert1991} (entanglement-based,
        using Bell-state measurements to detect eavesdropping) and
        \textbf{B92}~\cite{bennett1992} (two non-orthogonal states) within
        the same platform would enable direct performance and security
        comparisons with BB84~\cite{ieee10872935}. This is particularly
        relevant given that E91 relies on entanglement rather than
        key-sifting, providing a fundamentally different security guarantee.

  \item \textbf{Delayed measurement via quantum memory.}
        A recent proposal~\cite{arxiv2410} demonstrates that if Bob stores
        received qubits in a quantum memory and defers measurement until after
        Alice publicly announces her basis choices, the sifting step is
        eliminated entirely --- all transmitted qubits contribute to the
        final key, doubling the raw key rate compared to standard BB84.
        Simulating this variant within Sequre would illustrate the practical
        impact of quantum memory on QKD efficiency.

  \item \textbf{Extended attack models: Trojan-horse and PNS.}
        Beyond the intercept-resend model currently implemented, two
        additional attack strategies are relevant for a complete security
        analysis. In a \textit{Trojan-horse attack}~\cite{arxiv2502}, Eve
        injects bright probe pulses into Alice's or Bob's optical components
        and analyses the back-reflected light to infer internal settings
        without interacting with the quantum channel. In a
        \textit{Photon Number Splitting attack}, Eve exploits multi-photon
        pulses emitted by weak coherent sources to intercept one photon while
        forwarding the rest~\cite{arxiv2509}. Simulating both would
        demonstrate the security boundaries of BB84 under hardware-level
        imperfections.

\end{itemize}

\begin{thebibliography}{9}
  \bibitem{lo2005decoy}
    H.-K. Lo, X. Ma, K. Chen,
    \textit{Decoy State Quantum Key Distribution},
    Physical Review Letters, 94, 230504, 2005.

  \bibitem{arxiv2405}
    M. Zahidy et al.,
    \textit{Comparison of non-decoy single-photon source and decoy weak
    coherent pulse in quantum key distribution},
    arXiv:2405.19963, 2024.

  \bibitem{arxiv2505}
    arXiv:2505.20838 --- \textit{QKD channel attenuation modelling}, 2025.

  \bibitem{arxiv2503}
    \textit{High-rate self-referenced CV-QKD over high-loss free-space channel},
    arXiv:2503.10168, 2025.

  \bibitem{brassard1994}
    G. Brassard, L. Salvail,
    \textit{Secret-Key Reconciliation by Public Discussion},
    EUROCRYPT 1993, LNCS 765, pp.\ 410--423, 1994.

  \bibitem{ekert1991}
    A. K. Ekert,
    \textit{Quantum Cryptography Based on Bell's Theorem},
    Physical Review Letters, 67(6), pp.\ 661--663, 1991.

  \bibitem{bennett1992}
    C. H. Bennett,
    \textit{Quantum Cryptography Using Any Two Nonorthogonal States},
    Physical Review Letters, 68(21), pp.\ 3121--3124, 1992.

  \bibitem{ieee10872935}
    \textit{A Novel Approach Based on QKD Using BB84 and E91 Protocol for
    Resilient Encryption and Eavesdropper Detection},
    IEEE Xplore, DOI: 10.1109/ACCESS.2025.10872935.

  \bibitem{arxiv2410}
    \textit{Improving BB84 Efficiency with Delayed Measurement via Quantum
    Memory}, arXiv:2410.21191, 2024.

  \bibitem{arxiv2502}
    \textit{Imperfect Preparation and Trojan Attack on the Phase Modulator
    in the Decoy-State BB84 Protocol}, arXiv:2502.21160, 2025.

  \bibitem{arxiv2509}
    \textit{Resisting Quantum Key Distribution Attacks Using Quantum Machine
    Learning}, arXiv:2509.14282, 2025.
\end{thebibliography}

% ===== 5. CONCLUSIONS ===========================================
\section{Conclusions}

\textbf{Sequre} represents a complete, end-to-end implementation of the BB84
quantum key distribution protocol, with an emphasis on real-time visualization
and the accessibility of quantum cryptography concepts. The integration with IBM
Quantum hardware adds a practically unique dimension compared to purely simulated
solutions, enabling direct observation of real decoherence noise effects on key
exchange outcomes.

The project aims to demonstrate that quantum security is not an abstract concept
but a physically verifiable reality, and to serve as a pedagogically useful tool
for students and researchers interested in quantum cryptography.

Looking forward, Sequre is intended to grow beyond a thesis prototype into
a platform that actively \textbf{tests the boundaries of QKD} --- from
hardware-induced decoherence on IBM backends to advanced attack scenarios
such as PNS and Trojan-horse interception --- and that can find a
\textbf{practical use} in research, education, and as a building block in
future secure-communication systems where quantum-safe key exchange becomes
operationally required.

\vspace{24pt}
\noindent
\begin{tabular}{@{}p{7cm}p{7cm}@{}}
  \textbf{Student,}          & \textbf{Supervisor,} \\[24pt]
  Vezeteu Andrei             & Conf.\ Dr.\ Andreea Arusoaie \\
\end{tabular}

\end{document}
